% Part: methods
% Chapter: induction
% Section: structural-induction

\documentclass[../../../include/open-logic-section]{subfiles}

\begin{document}

\olfileid{mth}{ind}{sti}

\olsection{গঠনগত আরোহ}

এ পর্যন্ত আমরা আরোহ দিয়ে সব স্বাভাবিক সংখ্যা সম্পর্কে
ফল প্রতিষ্ঠা করেছি। কিন্তু অনুরূপ একটি নীতি আরোহীভাবে
সংজ্ঞায়িত সেটের সব !!{element}s সম্পর্কে সরাসরি ফল
প্রমাণ করতেও ব্যবহার করা যায়। একে প্রায়ই \emph{গঠনগত}
আরোহ বলে, কারণ এটি ওই বস্তুগুলির গঠনের উপর নির্ভর করে।

সাধারণত আরোহী সংজ্ঞায় থাকে (a)~সেটের ``প্রারম্ভিক''
!!{element}s-এর তালিকা এবং (b)~পুরোনো !!{element}s থেকে
নতুনগুলি তৈরির ক্রিয়াগুলির তালিকা। সুগঠিত পদের ক্ষেত্রে
প্রারম্ভিক বস্তুগুলি অক্ষর। এখানে ক্রিয়া মাত্র একটি:
\begin{align*}
  o(s_1, s_2) = & [s_1 \circ s_2]
\end{align*}
এমনকি স্বাভাবিক সংখ্যা~$\Nat$-কেও আরোহীভাবে সংজ্ঞায়িত
ভাবতে পারো: প্রারম্ভিক বস্তু~$0$, আর ক্রিয়াটি উত্তরসূরি
অপেক্ষক~$x + 1$।

আরোহীভাবে সংজ্ঞায়িত সেটের সব উপাদান সম্পর্কে কিছু প্রমাণ
করতে—অর্থাৎ প্রতিটি !!{element}-এর একটি ধর্ম~$P$
আছে দেখাতে—আমাদের করতে হবে:
\begin{enumerate}
\item প্রারম্ভিক বস্তুগুলির ধর্ম~$P$ আছে প্রমাণ।
\item প্রতিটি ক্রিয়া~$o$-এর ক্ষেত্রে, তার আর্গুমেন্টগুলির
  ধর্ম~$P$ থাকলে ফলেরও তা আছে প্রমাণ।
\end{enumerate}
যেমন, সব সুগঠিত পদ সম্পর্কে কিছু প্রমাণে প্রথমে সব
অক্ষরের জন্য তা দেখাই; তারপর $s_1$ ও $s_2$-এর
প্রত্যেকটির জন্য সত্য হলে $[s_1 \circ s_2]$-এর জন্যও
সত্য দেখাই।

\begin{prop}
  যেকোনো সুগঠিত পদ~$t$-তে $[$ ও $]$-এর সংখ্যা সমান।
\end{prop}

\begin{proof}
গঠনগত আরোহ ব্যবহার করি। সুগঠিত পদগুলি আরোহীভাবে
সংজ্ঞায়িত: অক্ষরগুলি প্রারম্ভিক বস্তু, আর পুরোনো পদ থেকে
নতুন পদ গড়ার ক্রিয়াটি~$o$।
\begin{enumerate}
\item প্রতিটি অক্ষরের জন্য দাবি সত্য: একক অক্ষরে $[$-এর
  সংখ্যা~$0$, আর $]$-এর সংখ্যাও~$0$।
\item ধরা যাক, $s_1$-এ $[$-এর সংখ্যা $]$-এর সংখ্যার
  সমান; $s_2$-এর ক্ষেত্রেও তাই। $o(s_1, s_2)$, অর্থাৎ
  $[s_1 \circ s_2]$-তে $[$-এর সংখ্যা হলো $s_1$ ও
  $s_2$-তে $[$-এর সংখ্যার যোগফলের চেয়ে এক বেশি।
  $o(s_1, s_2)$-তে $]$-এর সংখ্যাও $s_1$ ও $s_2$-তে
  $]$-এর সংখ্যার যোগফলের চেয়ে এক বেশি। তাই
  $o(s_1, s_2)$-তে $[$-এর সংখ্যা $o(s_1,s_2)$-তে
  $]$-এর সংখ্যার সমান।
\end{enumerate}
\end{proof}

\begin{prob}
  গঠনগত আরোহে প্রমাণ করো যে কোনো সুগঠিত পদ $]$
  দিয়ে শুরু হয় না।
\end{prob}

গঠনগত আরোহে আরেকটি প্রমাণ দেখা যাক। প্রতীক-স্ট্রিং~$t$-এর
একটি \emph{প্রকৃত প্রারম্ভিক খণ্ড} হলো এমন স্ট্রিং~$s$,
যার প্রতীকগুলি বাঁ দিক থেকে পড়লে~$t$-এর প্রতীকগুলির
সঙ্গে মেলে, কিন্তু $t$ আরও দীর্ঘ। যেমন, $[a \circ {}$
হলো $[a \circ b]$-এর প্রকৃত প্রারম্ভিক খণ্ড; কিন্তু
$[b \circ {}$ নয় (দ্বিতীয় প্রতীকই আলাদা), এবং
$[a \circ b]$-ও নয় (একই দৈর্ঘ্য)। নিচের দাবিতে
খালি স্ট্রিংটি বাদ দেওয়া হচ্ছে।

\begin{prop}\ollabel{prop:initial}
  % BN-SRC-485: the empty prefix is proper but has equal bracket counts.
  সুগঠিত পদ~$t$-এর প্রত্যেক অশূন্য প্রকৃত প্রারম্ভিক
  খণ্ডে $[$-এর সংখ্যা $]$-এর চেয়ে বেশি।
\end{prop}

\begin{proof}
~$t$-এর গঠনের উপর আরোহ করি:
  \begin{enumerate}
  \item $t$ একক অক্ষর: তাহলে $t$-এর কোনো অশূন্য প্রকৃত
    প্রারম্ভিক খণ্ড নেই।
  \item কোনো সুগঠিত পদ $s_1$ ও~$s_2$ দিয়ে
    $t = [s_1 \circ s_2]$। $r$ যদি $t$-এর অশূন্য
    প্রকৃত প্রারম্ভিক খণ্ড হয়, কয়েকটি সম্ভাবনা আছে:
    \begin{enumerate}
    \item $r$ কেবল $[$: $r$-এ $]$-এর তুলনায় একটি
      $[$ বেশি।
    \item $r$ হলো $[r_1$, যেখানে $r_1$ হলো~$s_1$-এর
      অশূন্য প্রকৃত প্রারম্ভিক খণ্ড: $s_1$ সুগঠিত পদ
      বলে আরোহের অনুমানে $r_1$-এ $]$-এর তুলনায় $[$
      বেশি; $[r_1$-এও তাই।
    \item $r$ হলো $[s_1$ অথবা $[s_1 \circ {}$:
      আগের ফল অনুসারে~$s_1$-এ $[$ ও $]$-এর সংখ্যা সমান।
      তাই $[s_1$ বা $[s_1 \circ {}$-তে $]$-এর চেয়ে
      একটি $[$ বেশি।
    \item $r$ হলো $[s_1 \circ r_2$, যেখানে $r_2$ হলো
      $s_2$-এর অশূন্য প্রকৃত প্রারম্ভিক খণ্ড:
      আরোহের অনুমানে $r_2$-এ $]$-এর চেয়ে $[$ বেশি।
      আগের ফল অনুসারে~$s_1$-এ $[$ ও $]$-এর সংখ্যা সমান।
      ফলে $[s_1 \circ r_2$-তে $]$-এর চেয়ে $[$ বেশি।
    \item $r$ হলো $[s_1 \circ s_2$: আগের ফল অনুসারে
      $s_1$ ও~$s_2$ উভয় পদে $[$ ও $]$-এর সংখ্যা সমান।
      তাই $[s_1 \circ s_2$-তে $]$-এর চেয়ে একটি $[$ বেশি।
    \end{enumerate}
  \end{enumerate}
\end{proof}

\end{document}
